\documentclass{article}
\usepackage[utf8]{inputenc}

\begin{document}


\tableofcontents


\section{Gruppedannelse}
Gruppen var selvvalgt til dette projekt, og det var forventet at man dannede grupper på 2-3 medlemmer (1 eller 4 personers grupper var også tilladt, men ikke anbefalet).
Vi valgte at danne denne gruppe, da vi alle går på samme semester og derigennem har fået et godt kendskab til hinanden. 
Derudover har vi også valgt at være i gruppe i vores andet fag det her semester (Web Arkitektur og Orkistrering), da det gav praktisk mening i forhold til planlægning af mødetider, og fordi det ville give en fleksibilitet da vi ville få muligheden for at skifte mellem fagene efter behov.


\section{Samarbejdsaftale}
Som vi har lært i tidligere semester er det en god ide at få skrevet en samarbejdsaftale til gruppeprojekter. 
Det har vi også valgt at gøre til dette projekt, og det var hermed en af de første ting vi fik lavet.
Den er meget standard, og man skal blandt andet give kage hvis man kommer forsent. [I bilag]


\section{Udviklingsforløb}
I dette projekt har vi valgt at afvige lidt fra udviklingsformer som vi har lært om i løbet af udannelsen.
Vi har valgt en meget agil arbejdsproces, der har været meget iterationsbaseret.
Det startede med at vi udvalgte og opstilte krav efter MoSCoW.
Herefter har vi løbende udvalgt issues som vi skulle udvikle som en del af et sprint.
Hvert sprint blev så opstillet i et Kanban-board, hvor vi så har kunne fordele diverse issues imellem os.
Som en afvigelse fra hvad vi har lært i undervisningen har vi som en anbefaling fra vores vejleder valgt at vente med at opstille en accepttest der tester hvert et krav fra vores MoSCoW.
I stedet har vi besluttet at udspecificer hvad testene tilhørende de features i sprintet skal indeholde.
Det vil sige at vi kun har opstillet test for de features vi har valgt at inkludere i produktet.
Grundlaget for dette valg er at vi hermed ikke skulle bruge tid på at opstille tests for features vi alligevel ikke vil have haft tid til udvikle.
\\
Vi startede ud med et relativt stort sprint hvor vi fik lavet alt det grundliggende for en funktionel MVP.
Her fik vi opsat vores database, backend og frontend, så vi havde en web-application med en login-funkionalitet, samt et database-schema der ville dække vores umidelbare behov for største delen af vores kommende features.
\\ 
Herefter opstilte vi lidt mindre sprint, hvor vi løbende udvidede vores web-apps funktionsdygtighed. 

\section{Projektadministration}
Til administrering af vores projekt har vi gjort brug af diverse platforme.
Der blev blandt andet brug GitHub til versionskontrol, samt deling af vores kode og rapport.
Vi besluttede os for at lave et repository som indeholdte alt frontend kode, backend kode, og rapporten. 
Dette besluttede vi fordi at vi havde så lille et hold, og det gjorde det nemmere at sikre at man altid havde versioner af frontend og backend som kunne spille sammen.
Udover versinoskontrol, har vi også brugt GitHubs Kanban-board som beskrevet tidligere.
\\
Til kommunikation har vi hovedsagelidt brugt Messenger og Discord. 
Messenger har vært brugt til daglig kommunikation og planlægning af møder. 
Hvorimod Discord har været brugt til deling af CLI-komandoer og code-snippets, og til deling af diverse resurser som PDF-filer.
\\
Til sidst har vi brugt Google Drive til deling af bl.a skitser og andre dokumenter der nemmere kunne hånteres der.
Vi førte også logbog i Google Drive.

\section{Individuelle konklusioner}

\subsection{Oliver}

\subsection{Thomas}

\subsection{Thor}

\end{document}